% =====================================================================
%  Artigo - Ferrugem asiatica da soja (formato do modelo do professor)
%  Modelo: A4, Times, duas colunas, titulo e resumo em uma coluna,
%  secoes "1 INTRODUCAO", subsecoes "2.1. Titulo" em italico,
%  legenda de figura embaixo ("Figura N - ...") e titulo de tabela
%  em cima ("Tabela N. ..."), referencias numeradas [1], [2]...
%
%  Linhas com % sao comentarios: nao aparecem no PDF.
%  Os topicos de cada secao vem do "Roteiro do Artigo".
% =====================================================================
\documentclass[10pt,a4paper,twocolumn]{article}

% --- Idioma e fonte ---
\usepackage[utf8]{inputenc}
\usepackage[T1]{fontenc}
\usepackage[brazilian]{babel}
\usepackage{newtxtext,newtxmath}   % Times, como no modelo

% --- Pagina ---
\usepackage[a4paper,top=2.2cm,bottom=2.2cm,left=2cm,right=2cm,columnsep=0.8cm]{geometry}
\raggedbottom
\setlength{\parindent}{0pt}
\setlength{\parskip}{0pt}

% --- Pacotes ---
\usepackage{graphicx}
\usepackage{amsmath}
\usepackage{booktabs}
\usepackage{float}
\usepackage{tikz}
\usetikzlibrary{positioning}
\usepackage[hidelinks]{hyperref}

% --- Titulos de secao no estilo do modelo ---
\usepackage{titlesec}
\titleformat{\section}{\normalfont\normalsize\MakeUppercase}{\thesection}{0.6em}{}
\titleformat{\subsection}{\normalfont\normalsize}{\thesubsection.}{0.4em}{\itshape}
\titlespacing*{\section}{0pt}{10pt}{6pt}
\titlespacing*{\subsection}{0pt}{8pt}{4pt}

% --- Legendas: "Figura 1 – ..." embaixo; "Tabela 1. ..." em cima ---
\usepackage{caption}
\captionsetup{font=small}
\captionsetup[figure]{labelsep=endash}
\captionsetup[table]{labelsep=period,position=top}

% --- Figura que compila mesmo antes de enviar a imagem ao Overleaf ---
% Uso: \figura{largura}{arquivo}
\newcommand{\figura}[2]{%
  \IfFileExists{#2}{\includegraphics[width=#1]{#2}}{%
    \fbox{\parbox[c][3cm][c]{\dimexpr #1-2\fboxsep-2\fboxrule\relax}{\centering\footnotesize
      Enviar ao Overleaf:\\\texttt{\detokenize{#2}}}}}}

% --- Nome da lista de referencias ---
\addto\captionsbrazilian{\renewcommand{\refname}{REFERÊNCIAS}}

\begin{document}

% =====================================================================
%  TITULO, AUTORES E RESUMO (uma coluna)
% =====================================================================
\twocolumn[{%
\begin{center}
{\Large Quantificação de clorose e necrose da ferrugem asiática da soja por segmentação em superpixels e atributos no espaço CIELab\par}
\end{center}
\vspace{6pt}
Heitor Bartolomei Pauleto e \textit{Pedro Hiniralow Canato}\\
\textit{Curso de Ciência da Computação, Universidade Paulista (UNIP), São José do Rio Preto -- SP, Brasil.}

\vspace{10pt}
RESUMO: A ferrugem asiática da soja, causada pelo fungo
\textit{Phakopsora pachyrhizi}, provoca clorose e necrose nas folhas, e
a proporção dessas regiões indica a severidade da doença. Este trabalho
apresenta um método de processamento digital de imagens para estimar
essa severidade a partir de imagens RGB, sem aprendizado profundo. Cada
folíolo é dividido em superpixels pelo algoritmo SLIC, e cada superpixel
é classificado como saudável, clorótico ou necrótico por regras sobre o
matiz e a luminosidade no espaço CIELab. O método foi aplicado a 68
imagens de 18 folíolos de um conjunto público, e os limiares foram
calibrados contra anotação manual, com divisão entre treino e teste por
folíolo. Nas 22 imagens anotadas, a severidade estimada teve correlação
de 0,87 com a anotada, e o F1 foi de 0,82 para a necrose e 0,63 para a
clorose, que se mostrou a principal fonte de erro. A comparação entre
três algoritmos de superpixel e a classificação pixel a pixel, em CIELab
e HSV, indicou impacto pequeno na precisão, e o Watershed compacto foi
uma alternativa equivalente e mais rápida ao SLIC. A análise temporal
confirmou o aumento da severidade ao longo dos dias e revelou
inconsistências na nomenclatura das imagens do conjunto público.

\vspace{4pt}
Palavras-chave: ferrugem asiática da soja; superpixel; CIELab; clorose; necrose; processamento digital de imagens.

\vspace{8pt}
ABSTRACT: Asian soybean rust, caused by the fungus \textit{Phakopsora
pachyrhizi}, produces chlorosis and necrosis on the leaves, and the
proportion of these regions indicates disease severity. This work
presents a digital image processing method to estimate this severity
from RGB images, without deep learning. Each leaflet is divided into
superpixels by the SLIC algorithm, and each superpixel is classified as
healthy, chlorotic or necrotic by rules on hue and lightness in the
CIELab color space. The method was applied to 68 images of 18 leaflets
from a public dataset, and the thresholds were calibrated against manual
annotation, with a leaflet-wise split between training and test sets. On
the 22 annotated images, the estimated severity had a correlation of
0.87 with the annotated one, and the F1-score was 0.82 for necrosis and
0.63 for chlorosis, which proved to be the main source of error. A
comparison of three superpixel algorithms and pixel-wise classification,
in CIELab and HSV, showed little impact on accuracy, and the compact
Watershed was an equivalent and faster alternative to SLIC. The temporal
analysis confirmed the increase in severity over the days and revealed
naming inconsistencies in the images of the public dataset.

\vspace{4pt}
Keywords: Asian soybean rust; superpixel; CIELab; chlorosis; necrosis; digital image processing.
\vspace{14pt}
}]

% =====================================================================
\section{Introdução}

A soja tem grande importância na agricultura mundial, e o Brasil
responde por cerca de um terço da produção global \cite{dasilva2022}.
Entre as doenças que afetam a cultura está a ferrugem asiática, causada
pelo fungo \textit{Phakopsora pachyrhizi}. A infecção provoca clorose
(amarelecimento) e necrose (morte do tecido) nas folhas; a proporção das
áreas afetadas é um indicador da severidade da doença.

A avaliação visual da severidade pode ser demorada e subjetiva. Como
alternativa, o processamento digital de imagens (PDI) permite analisar
fotografias RGB de forma automática e reprodutível. Da Silva et al.
\cite{dasilva2022} exploraram essa abordagem usando o espaço de cor
CIELab e superpixels.

Este trabalho implementa e amplia essa abordagem sem recorrer ao
aprendizado profundo. Além da validação com anotações manuais, o estudo
calibra os limiares em um conjunto de treino separado, compara três
métodos de superpixel com a classificação pixel a pixel nos espaços
CIELab e HSV, e acompanha a severidade de cada folíolo pelas datas reais
das fotos. Essa análise temporal revelou divergências entre as datas das
capturas e a sequência sugerida pelos nomes dos arquivos.

O restante do artigo está organizado da seguinte forma: a seção 2 reúne
os trabalhos relacionados, a seção 3 descreve o método, a seção 4
apresenta os resultados e a seção 5 encerra com as conclusões.

% =====================================================================
\section{Trabalhos relacionados}
% Um paragrafo (ou subsecao 2.1, 2.2...) por trabalho:
% - da Silva et al. \cite{dasilva2022}: superpixel + intervalos CIELab;
%   indices de necrose, clorose e ferrugem.
% - Phan et al. \cite{phan2022}: SLIC + deep learning em milho; o
%   superpixel como etapa previa.
% - Khan e Narvekar \cite{khan2022}: balanceamento de cor + superpixel em
%   tomate. Liga com o trabalho futuro.
% - Feng et al. \cite{feng2024}: hiperespectral + deep learning, mais
%   preciso porem caro. Contrasta com a proposta (RGB, sem treino).
% - Lopes e Fernandes \cite{lopes2026}: CNN leve para lesoes de ferrugem.
%   So temos o resumo: citar apenas o que esta nele.

Da Silva et al. \cite{dasilva2022} segmentaram folíolos de soja em
regiões saudáveis, cloróticas e necróticas a partir do espaço CIELab,
utilizando superpixels como pós-processamento. O método gera índices de
clorose, necrose e ferrugem e é a base deste trabalho.

Phan et al. \cite{phan2022} dividiram folhas de milho em superpixels
com o SLIC e classificaram cada região com redes convolucionais
pré-treinadas, chegando a 97,77\% de acurácia em imagens de campo.

Khan e Narvekar \cite{khan2022} combinaram balanceamento de cores e
superpixels para detectar doenças em tomate em ambiente natural, com
acurácia de cerca de 93\%. O balanceamento de cores é tratado aqui como
trabalho futuro.

Feng et al. \cite{feng2024} detectaram a ferrugem asiática da soja antes
dos sintomas com imagens hiperespectrais e uma rede convolucional 3D,
com acurácia próxima de 97\%, mas o método exige equipamento caro e
treinamento profundo.

Lopes e Fernandes \cite{lopes2026} propuseram o ASRDNet, uma rede leve
para segmentar lesões de ferrugem em soja, com F1 de 0,9251 e menos
parâmetros que U-Net, SegNet e DeepLabv3.

Diferentemente desses trabalhos com aprendizado profundo, este artigo
usa apenas imagens RGB e regras de cor, sem treinamento, e avalia o
método contra anotação manual.

% =====================================================================
\section{Método}
% Frase de abertura: o metodo classifica cada superpixel da folha como
% saudavel, clorose ou necrose pela cor em CIELab, com limiares
% calibrados contra anotacao manual. A Figura 1 resume as etapas.

O método classifica cada superpixel da folha como saudável, clorose ou necrose pela cor no espaço CIELab, com limiares calibrados contra anotação manual. A Figura~\ref{fig:fluxo} resume as etapas, que são descritas nas subseções a seguir.

\begin{figure}[H]
\centering
\begin{tikzpicture}[
  node distance=3.5mm,
  box/.style={draw,rounded corners=2pt,align=center,font=\footnotesize,
              minimum width=0.8\columnwidth,inner sep=3pt},
  every path/.style={->,>=stealth}]
  \node[box] (a) {Imagem RGB do folíolo};
  \node[box,below=of a] (b) {Redução ($\leq$ 900 px) e máscara da folha\\(limiar de branco, abertura e erosão)};
  \node[box,below=of b] (c) {SLIC ($\approx$ 4.000 superpixels, compacidade 10)};
  \node[box,below=of c] (d) {Média de $L^*a^*b^*$ por superpixel;\\matiz $h$ e $L^*$};
  \node[box,below=of d] (e) {Regras: necrose se $h<75^\circ$ ou $L^*<55$;\\clorose se $h<96^\circ$; senão saudável};
  \node[box,below=of e] (f) {Severidade (\%) = (clorose + necrose) / folha};
  \draw (a)--(b); \draw (b)--(c); \draw (c)--(d); \draw (d)--(e); \draw (e)--(f);
\end{tikzpicture}
\caption{Fluxograma do método proposto.}
\label{fig:fluxo}
\end{figure}

\subsection{Dados}

Foram utilizadas 68 imagens de 18 folíolos de soja destacados,
correspondentes ao Experimento 2 de da Silva et al. \cite{dasilva2022} e
disponíveis no repositório público dos autores\footnote{\url{https://github.com/diego91964/paper_compag_106542}}.
As imagens já são fornecidas com o fundo branco removido pelos autores.

Cada folíolo foi fotografado em várias sessões ao longo da infecção, e
os arquivos são nomeados no formato \textit{folíolo\_captura} (por
exemplo, 016\_3 é a terceira captura do folíolo 016). As datas de cada
foto foram lidas dos metadados EXIF, o que permitiu organizar as imagens
em seis sessões, nos dias 0, 3, 6, 9, 14 e 17 após a primeira captura.
Essa ordenação pelas datas reais é a base da análise temporal
apresentada na seção 4.

\subsection{Pré-processamento e superpixels}

Para reduzir o tempo de processamento, cada imagem foi redimensionada
para no máximo 900 pixels no maior lado. Em seguida, a máscara da folha
foi obtida por um limiar de branco: um pixel é considerado parte da
folha quando $\min(R,G,B) < 235$, o que separa o tecido vegetal do fundo
branco. Pequenas falhas dessa máscara foram corrigidas com uma abertura
($5\times5$) e um fechamento ($9\times9$). Em seguida, uma abertura com
elemento estruturante elíptico de raio igual a 2,2\% do maior lado da
imagem remove o pecíolo, que é estreito, e as bordas arredondadas da
folha eliminadas por essa operação são recuperadas quando estão próximas
do corpo da folha. Mantém-se apenas o maior componente conexo, e uma
erosão com raio de 0,7\% do maior lado descarta a borda, onde há mistura
com o fundo e sombra do recorte.

A região da folha foi então dividida em superpixels pelo algoritmo SLIC
\cite{achanta2012}, com cerca de 4.000 segmentos e compacidade 10. A
segmentação foi aplicada apenas dentro da máscara, de modo que o fundo
não gera superpixels. Cada superpixel agrupa pixels vizinhos de cor
semelhante e passa a ser tratado como uma unidade na etapa de
classificação. A implementação foi feita em Python, com as bibliotecas
scikit-image \cite{skimage2014} e OpenCV \cite{opencv2000}, e o código
está disponível publicamente\footnote{\url{https://github.com/Heitor117/deteccao_de_ferrugem_na_soja}}.

\subsection{Classificação e severidade}

Para cada superpixel, calculou-se a média dos componentes $L^*$, $a^*$ e
$b^*$ dos seus pixels, com $a^*$ e $b^*$ centralizados em zero (o OpenCV
os armazena com deslocamento de 128). A partir deles obteve-se o ângulo
de matiz:
\begin{equation}
h = \operatorname{atan2}(b^*, a^*)
\end{equation}
que indica a cor do tecido: valores maiores correspondem ao verde da
folha saudável, e valores menores, ao amarelo e ao marrom das lesões. A
luminosidade $L^*$, na escala do OpenCV (0 a 255), complementa o matiz
na identificação do tecido necrosado, que é mais escuro.

Cada superpixel foi classificado pela seguinte regra:
\begin{equation}
\text{classe} =
\begin{cases}
\text{necrose}, & h < 75^\circ \text{ ou } L^* < 55\\
\text{clorose}, & h < 96^\circ\\
\text{saudável}, & \text{caso contrário}
\end{cases}
\end{equation}
As condições são verificadas nessa ordem. Superpixels muito claros e
pouco saturados ($L^* > 185$ e croma $C^* = \sqrt{a^{*2} + b^{*2}} < 14$),
correspondentes a reflexos de luz, foram ignorados e não entram na área
da folha. Os limiares foram calibrados contra
anotação manual, como descrito na seção 3.4.

Por fim, a severidade da folha foi calculada como a porcentagem da área
afetada:
\begin{equation}
S = \frac{A_{\text{clorose}} + A_{\text{necrose}}}{A_{\text{folha}}}
\times 100
\end{equation}
em que $A$ é a área, em pixels, de cada classe. A Figura~\ref{fig:exemplo}
mostra o resultado em um folíolo do conjunto de teste.

\begin{figure}[H]
\centering
\figura{\columnwidth}{ImagensParaUtilizacao/img/016_3.png}
\caption{Folíolo 016 na captura 3: imagem original (esquerda), divisão em
3.969 superpixels pelo SLIC (centro) e classificação em saudável (verde),
clorose (laranja) e necrose (marrom escuro), com severidade de 11,93\%
(direita).}
\label{fig:exemplo}
\end{figure}

\subsection{Anotação e calibração}

Para avaliar o método, 22 imagens foram anotadas manualmente na
ferramenta Labelme \cite{labelme}, com polígonos delimitando as regiões
de clorose e de necrose. O restante da folha foi considerado saudável.
Foram anotadas apenas as lesões visíveis, e pontos muito pequenos foram
ignorados.

As imagens anotadas foram divididas por folíolo, para que capturas da
mesma folha não ficassem ao mesmo tempo no treino e no teste. O conjunto
de teste reuniu os folíolos 001, 012 e 016 (8 imagens), e as 14 imagens
restantes formaram o conjunto de treino.

Os limiares iniciais foram ajustados visualmente: $100^\circ$ para o
matiz da clorose, $80^\circ$ para o matiz da necrose e 55 para a
luminosidade da necrose, com uma regra adicional de contraste local:
um superpixel cujo matiz ficasse mais de $12^\circ$ abaixo da mediana dos
vizinhos, num raio de 60 pixels, era tratado como lesão. Esses valores
foram então calibrados por busca em grade no conjunto de treino,
resultando em $96^\circ$, $75^\circ$ e 55. A regra de contraste local
foi descartada, pois piorava o resultado. Os limiares
calibrados foram aplicados, sem novos ajustes, ao conjunto de teste.

\subsection{Métricas e análise temporal}

A classificação foi comparada à anotação manual pixel a pixel. Para cada
classe (saudável, clorose e necrose), foram calculados o F1 e o IoU
(\textit{Intersection over Union}):
\begin{equation}
F1 = \frac{2\,VP}{2\,VP + FP + FN}, \qquad
IoU = \frac{VP}{VP + FP + FN}
\end{equation}
em que $VP$, $FP$ e $FN$ são os pixels verdadeiros positivos, falsos
positivos e falsos negativos da classe. A severidade estimada foi
comparada à anotada pelo erro absoluto médio, em pontos percentuais
(p.p.), e pelo coeficiente de correlação de Pearson ($r$).

Na análise temporal, a severidade de cada folíolo foi acompanhada ao
longo dos dias de captura. A progressão da doença foi resumida pela
AUDPC (área sob a curva de progresso da doença) \cite{shaner1977},
calculada pela regra do trapézio:
\begin{equation}
\text{AUDPC} = \sum_{i=1}^{n-1} \frac{S_i + S_{i+1}}{2}\,(t_{i+1} - t_i)
\end{equation}
em que $S_i$ é a severidade na sessão $i$ e $t_i$ é o dia
correspondente. Como o número de sessões varia entre os folíolos, a
AUDPC também foi dividida pela duração do acompanhamento
($t_n - t_1$), resultando na AUDPC padronizada (sAUDPC), que equivale à
severidade média no período. Quando o mesmo folíolo foi fotografado
duas vezes na mesma sessão, usou-se a média das duas fotos.

\subsection{Comparação de segmentações}

Para avaliar o papel do superpixel e do espaço de cor, foram comparadas
quatro formas de segmentar a folha: SLIC \cite{achanta2012},
Felzenszwalb \cite{felzenszwalb2004}, Watershed compacto
\cite{neubert2014} e a classificação pixel a pixel, sem superpixel. Os
três algoritmos foram configurados para gerar cerca de 4.000 regiões por
folha: o Felzenszwalb com escala 10, $\sigma = 0{,}5$ e tamanho mínimo de
20 pixels, e o Watershed compacto sobre o gradiente de Sobel da imagem em
tons de cinza, com 4.000 marcadores e compacidade 0,001. Cada forma de
segmentação foi testada em CIELab e em HSV, totalizando oito
configurações. No HSV, a mesma regra usa o canal H como matiz e o canal
V como brilho, e a saturação S substitui o croma na exclusão de
reflexos. Para cada configuração, os três limiares foram recalibrados
por busca em grade no conjunto de treino, escolhendo a combinação de
maior F1 médio, e então avaliados no conjunto de teste. A Figura~\ref{fig:segmentacoes} mostra as
fronteiras obtidas por cada método.

\begin{figure}[H]
\centering
\figura{\columnwidth}{ImagensParaUtilizacao/comparacao/fig_segmentacoes.png}
\caption{Fronteiras obtidas pelos métodos de segmentação comparados.}
\label{fig:segmentacoes}
\end{figure}

% =====================================================================
\section{Testes e resultados}
% Abertura: o metodo acompanha bem a anotacao (r = 0,87) e detecta bem a
% necrose; a clorose e o ponto fraco, e a escolha do superpixel quase
% nao muda a precisao.

O método acompanha bem a anotação manual, com correlação de Pearson $r = 0{,}87$ entre a severidade anotada e a estimada nas 22 imagens anotadas, e detecta bem a necrose. A clorose é o ponto fraco, e a escolha do algoritmo de superpixel quase não altera a precisão. As subseções a seguir detalham a calibração, a concordância com a anotação, a evolução da doença no tempo e a comparação entre segmentações.

\subsection{Calibração}
% A calibracao generaliza: no teste, F1 medio 0,64 -> 0,68 e erro
% 20,1 -> 17,9 pontos (Tabela 1).

A calibração melhorou o desempenho no treino e manteve o ganho nas imagens de teste (Tabela~\ref{tab:calibracao}). No treino, o F1 médio passou de 0,70 para 0,73 e o erro da severidade caiu de 17,1 para 12,4 pontos percentuais. No teste, o F1 médio passou de 0,65 para 0,69 e o erro caiu de 20,1 para 17,9 pontos. No teste, o ganho aparece nas três classes: o F1 foi de 0,75 para 0,79 no tecido saudável, de 0,53 para 0,54 na clorose e de 0,67 para 0,73 na necrose, que foi a classe mais beneficiada.

\begin{table}[H]
\centering
\caption{Limiares padrão (100°/80°/55, com contraste local) e calibrados (96°/75°/55) no treino (14 imagens) e no teste (8 imagens).}
\label{tab:calibracao}
\small
\setlength{\tabcolsep}{3.5pt}
\begin{tabular}{llccccc}
\toprule
 & & \multicolumn{4}{c}{F1} & Erro \\
\cmidrule(lr){3-6}
Conjunto & Limiares & Saud. & Clor. & Necr. & Média & (p.p.) \\
\midrule
Treino & Padrão    & 0,63 & 0,66 & 0,81 & 0,70 & 17,1 \\
Treino & Calibrado & 0,68 & 0,67 & 0,84 & 0,73 & 12,4 \\
\midrule
Teste  & Padrão    & 0,75 & 0,53 & 0,67 & 0,65 & 20,1 \\
Teste  & Calibrado & 0,79 & 0,54 & 0,73 & 0,69 & 17,9 \\
\bottomrule
\end{tabular}
\end{table}

\subsection{Concordância com a anotação}
% Nas 22 anotadas (inclui o treino, entao e otimista): F1 saudavel 0,74,
% clorose 0,63, necrose 0,82; erro 14,4 pontos; r = 0,87.
% Principal erro: saudavel virando clorose (pegar o % da matriz).
% Maioria dos pontos acima da diagonal: tendencia a superestimar,
% sobretudo em folhas intermediarias com amarelado difuso.

Nas 22 imagens anotadas, o F1 foi de 0,74 para o tecido saudável, 0,63 para a clorose e 0,82 para a necrose, com erro médio de 14,4 pontos percentuais e $r = 0{,}87$ (Tabela~\ref{tab:avaliacao}). Como esse conjunto inclui as imagens de treino, os valores são otimistas e devem ser lidos junto com os resultados do conjunto de teste.

A necrose é a classe mais bem detectada, e a clorose é a mais difícil. A matriz de confusão (Figura~\ref{fig:matriz}), normalizada pela classe anotada, mostra que o principal erro é tecido saudável classificado como clorose, que atinge 28\% dos pixels saudáveis. Também há confusão entre as classes de lesão: 22\% dos pixels de clorose foram classificados como saudáveis, e 15\% dos pixels de necrose, como clorose. O primeiro erro se reflete na severidade: na Figura~\ref{fig:dispersao}, a maioria dos pontos fica acima da diagonal, ou seja, o método tende a superestimar a severidade, sobretudo em folhas intermediárias com amarelado difuso, em que a fronteira entre verde e amarelo é pouco nítida.

\begin{table}[H]
\centering
\caption{Desempenho por classe nas 22 imagens anotadas.}
\label{tab:avaliacao}
\small
\begin{tabular}{lcc}
\toprule
Classe & F1 & IoU \\
\midrule
Saudável & 0,74 & 0,58 \\
Clorose  & 0,63 & 0,46 \\
Necrose  & 0,82 & 0,69 \\
\midrule
\multicolumn{3}{l}{Erro da severidade: 14,4 p.p.; $r$ = 0,87} \\
\bottomrule
\end{tabular}
\end{table}

\begin{figure}[H]
\centering
\figura{0.85\columnwidth}{ImagensParaUtilizacao/figuras/fig_dispersao_anotado.png}
\caption{Severidade anotada versus severidade estimada.}
\label{fig:dispersao}
\end{figure}

\begin{figure}[H]
\centering
\figura{0.85\columnwidth}{ImagensParaUtilizacao/figuras/fig_matriz_confusao.png}
\caption{Matriz de confusão por pixel nas 22 imagens anotadas, normalizada pela classe anotada (cada linha soma 100\%).}
\label{fig:matriz}
\end{figure}

\begin{table*}[t]
\centering
\caption{Comparação das segmentações e espaços de cor no conjunto de teste (calibrados no treino).}
\label{tab:comparacao}
\small
\begin{tabular}{llrcccccc}
\toprule
Segmentação & Cor & Segmentos & Tempo de segm. (s) & F1 clorose & F1 necrose & F1 médio & Erro (p.p.) \\
\midrule
Watershed compacto & CIELab & 4.050   & 0,34  & 0,546 & 0,730 & 0,688 & 18,2 \\
SLIC               & CIELab & 3.946   & 20,54 & 0,540 & 0,731 & 0,685 & 17,9 \\
Felzenszwalb       & CIELab & 4.082   & 1,96  & 0,540 & 0,726 & 0,684 & 18,1 \\
Sem superpixel     & CIELab & 288.902 & 0,43  & 0,534 & 0,716 & 0,678 & 17,2 \\
Watershed compacto & HSV    & 4.050   & 0,34  & 0,540 & 0,717 & 0,678 & 18,7 \\
Felzenszwalb       & HSV    & 4.082   & 1,96  & 0,536 & 0,718 & 0,677 & 18,5 \\
SLIC               & HSV    & 3.946   & 20,54 & 0,537 & 0,717 & 0,677 & 18,2 \\
Sem superpixel     & HSV    & 288.902 & 0,43  & 0,532 & 0,713 & 0,674 & 17,5 \\
\bottomrule
\end{tabular}
\end{table*}

\subsection{Evolução no tempo}
% 68 imagens: severidade media 39,8%, mediana 26,1%, de 0,04% (010_2)
% a 100% (008_2). Os 18 foliolos tem severidade crescente com os dias.
% Foliolos 010 e 011: numero do arquivo nao segue a ordem cronologica
% (confirmado pelo EXIF); corrigido, a queda aparente some.
% Fotos do mesmo dia: 002_3 x 002_4 = 17,4 pontos; 004_3 x 004_4 = 2,8.

Nas 68 imagens, a severidade estimada tem média de 39,8\% e mediana de 26,1\%, variando de 0,04\% (imagem 010\_2) a 100\% (imagem 008\_2). Por sessão, a severidade média subiu de 17,3\% no dia 0 para 83,4\% no dia 17 (Tabela~\ref{tab:sessoes}), embora o número de folíolos fotografados diminua nas últimas sessões. Dos 16 folíolos fotografados em mais de uma sessão, 15 tiveram aumento de severidade (Figura~\ref{fig:barras}). A exceção foi o folíolo 010, que permaneceu praticamente sem sintomas, com severidade máxima de 2,45\%. Os folíolos 013 e 018 têm uma única captura e não permitem avaliar a evolução.

\begin{table}[H]
\centering
\caption{Severidade estimada por sessão de fotos (uma linha por folíolo e sessão).}
\label{tab:sessoes}
\small
\begin{tabular}{cccc}
\toprule
Dia & Folíolos & Média (\%) & Mediana (\%) \\
\midrule
0  & 17 & 17,3 & 4,7  \\
3  & 16 & 27,6 & 15,1 \\
6  & 14 & 40,3 & 32,1 \\
9  & 12 & 63,1 & 71,4 \\
14 & 4  & 67,6 & 70,1 \\
17 & 3  & 83,4 & 89,5 \\
\bottomrule
\end{tabular}
\end{table}

Como os folíolos foram acompanhados por períodos diferentes, de 2 a 6 sessões, a progressão foi comparada pela sAUDPC. Ela variou de 1,0\% no folíolo 010 a 97,4\% no folíolo 008, que já estava muito afetado na primeira captura. Entre os três folíolos acompanhados durante os 17 dias (006, 016 e 017), a AUDPC foi de 824,6, 460,9 e 305,1, respectivamente, o que mostra ritmos de progressão bem diferentes mesmo com o mesmo período de observação.

A ordenação pelas datas do EXIF revelou uma inconsistência no conjunto público: nos folíolos 010 e 011, o número no nome do arquivo não segue a ordem cronológica das fotos. Com as capturas reordenadas pela data real, a queda de severidade que aparecia nesses folíolos desaparece.

A severidade também varia entre fotos do mesmo dia. Entre as capturas 002\_3 e 002\_4 a diferença foi de 17,4 pontos percentuais, e entre 004\_3 e 004\_4 foi de 2,8. Isso sugere que a posição e a iluminação da folha em cada foto influenciam o resultado, além da evolução da doença.

\begin{figure}[H]
\centering
\figura{\columnwidth}{ImagensParaUtilizacao/figuras/fig_barras_imagens.png}
\caption{Evolução da severidade de um folíolo e severidade de todas as imagens.}
\label{fig:barras}
\end{figure}

\subsection{Comparação de segmentações e espaços de cor}
% CIELab melhor que HSV nas 4 segmentacoes (0,006 a 0,010 no F1).
% Superpixel reduz ~290 mil pixels a ~4 mil regioes e limpa o mapa,
% sem aumentar a precisao. Gargalo: clorose (F1 0,53 a 0,55 em todas).
% Diferencas de 0,01 com 8 imagens nao permitem dizer que um e melhor.
% Watershed compacto: equivalente e ~60x mais rapido que o SLIC.
% Tempos medidos com processos em paralelo: so para comparar.

A Tabela~\ref{tab:comparacao} e a Figura~\ref{fig:comparacao} comparam as oito configurações no conjunto de teste. O CIELab superou o HSV nas quatro segmentações, mas por pouco: a diferença no F1 médio foi de 0,004 a 0,010. A vantagem é pequena, porém consistente.

O superpixel reduz cerca de 290 mil pixels por folha a cerca de 4 mil regiões e limpa o mapa de classificação, mas não aumenta a precisão: em CIELab, o F1 médio ficou em 0,678 sem superpixel e entre 0,684 e 0,688 com ele, e o erro da severidade foi até menor sem superpixel (17,2 contra 17,9 a 18,2 pontos). O gargalo está na clorose, cujo F1 ficou entre 0,53 e 0,55 em todas as configurações, qualquer que seja a segmentação ou o espaço de cor.

Com apenas 8 imagens de teste, diferenças de 0,01 no F1 não permitem afirmar que um algoritmo é melhor que outro. Por isso o método principal continua sendo SLIC com CIELab, e o Watershed compacto entra como alternativa de desempenho equivalente e cerca de 60 vezes mais rápida na segmentação (0,34~s contra 20,54~s por folha). Os tempos medem apenas a etapa de segmentação, foram obtidos com vários processos em paralelo e servem somente para comparar os métodos, não como tempo absoluto. Na recalibração de cada configuração, a busca escolheu em todas as segmentações em CIELab os mesmos limiares do método principal ($96^\circ$, $75^\circ$ e 55), e $52^\circ$ e $36^\circ$ para os matizes no HSV.

\begin{figure}[H]
\centering
\figura{\columnwidth}{ImagensParaUtilizacao/comparacao/fig_comparacao.png}
\caption{Desempenho das oito configurações comparadas.}
\label{fig:comparacao}
\end{figure}

% --- Itens opcionais (use se sobrar espaco) ---
% figuras/fig_ordenado.png        - imagens ordenadas por severidade
% figuras/fig_evolucao_folhas.png - curvas de evolucao por foliolo
% figuras/fig_hist_matiz.png      - histograma do matiz com limiares
% figuras/fig_calibracao.png      - efeito do limiar de clorose
% tabelas: sessoes (dia, n, severidade media) e AUDPC por foliolo

% =====================================================================

\subsection{Discussão}

A principal fonte de erro do método está na clorose, cujo F1 ficou entre
0,53 e 0,55 em todas as configurações testadas. Parte dessa dificuldade
vem da própria natureza do sintoma: o amarelecimento se espalha de forma
gradual, sem uma borda definida entre o tecido verde e o clorótico.
Durante a anotação manual, essa transição também foi difícil de
delimitar, e os polígonos foram traçados apenas nas regiões de
amarelecimento evidente. Assim, áreas de amarelado leve, que o método
classifica como clorose pelo matiz, ficam marcadas como saudáveis na
anotação. Isso reduz o F1 da clorose e explica parte da superestimação
da severidade observada na Figura~\ref{fig:dispersao}. Ou seja, parte do
erro reflete a incerteza da própria referência, e não apenas uma falha
do método. Esse resultado reforça a necessidade de mais anotadores, que
permitiriam medir a concordância entre avaliadores e estimar o limite de
precisão alcançável para essa classe. 

\section{Conclusão}

Este trabalho apresentou um método para estimar a severidade da
ferrugem asiática da soja a partir de imagens RGB, combinando
superpixels SLIC com regras sobre o matiz e a luminosidade no espaço
CIELab, sem treinamento de redes neurais. Com limiares calibrados contra
anotação manual, o método acompanhou a severidade anotada ($r = 0{,}87$
nas 22 imagens anotadas) e detectou bem a necrose. O ponto fraco é a
clorose, principalmente em folhas com amarelado difuso, em que o método
tende a superestimar a área afetada.

A comparação entre segmentações mostrou que o algoritmo de superpixel e
o espaço de cor têm impacto pequeno na precisão. O superpixel simplifica
o mapa de classificação, mas não melhora o F1, e o CIELab teve vantagem
pequena, porém consistente, sobre o HSV. O Watershed compacto obteve
desempenho equivalente ao SLIC em muito menos tempo, sendo uma
alternativa viável. Assim, o limite do método está na separação entre
tecido saudável e clorótico, e não na etapa de segmentação. A análise
temporal, por sua vez, confirmou o aumento da severidade ao longo dos
dias e revelou inconsistências na nomenclatura do conjunto público, que
só puderam ser corrigidas com as datas reais das fotos.

O trabalho tem limitações. As anotações foram feitas por um único
anotador, com critério conservador, e apenas 22 imagens foram anotadas,
das quais 8 compõem o teste. Folíolos fotografados juntos, como 016 e
017, ficaram em conjuntos diferentes, o que cria uma dependência leve
entre treino e teste. Além disso, os limiares são globais e foram
definidos em um único experimento, com iluminação controlada.

Como trabalhos futuros, propõe-se aplicar balanceamento de cores antes
da classificação \cite{khan2022}, incluir atributos de textura e um
classificador aprendido, ampliar a anotação com mais avaliadores e
testar o método nas imagens de campo do mesmo repositório.

% =====================================================================
\section*{Agradecimentos}

Os autores agradecem a da Silva et al. \cite{dasilva2022} por disponibilizarem publicamente as imagens utilizadas neste trabalho e ao professor orientador pelo acompanhamento e pelas sugestões ao longo do desenvolvimento do artigo.

{\small
\bibliographystyle{unsrt}
\bibliography{refs}
}

\end{document}